% Generated draft: requires mathematical and source review.
\documentclass[11pt]{article}
\usepackage[margin=25mm]{geometry}
\usepackage{fontspec}
\setmainfont{Latin Modern Roman}
\usepackage{amsmath,amssymb,mathtools,unicode-math}
\setmathfont{Latin Modern Math}
\usepackage{enumitem,needspace,xurl,hyperref}
\setlength{\parindent}{0pt}
\setlength{\parskip}{5pt}
\setlength{\emergencystretch}{3em}
\newcommand{\N}{\mathbb{N}}
\newcommand{\Z}{\mathbb{Z}}
\newcommand{\Q}{\mathbb{Q}}
\newcommand{\R}{\mathbb{R}}
\newcommand{\C}{\mathbb{C}}
\newcommand{\F}{\mathbb{F}}
\newcommand{\A}{\mathbb{A}}
\newcommand{\PP}{\mathbb{P}}
\newcommand{\E}{\mathbb{E}}
\newcommand{\eps}{\varepsilon}
\newcommand{\dd}{\,\mathrm{d}}
\begin{document}
\section*{UnsolvedMath ID: 30000825}
\textbf{Source problem code:} catalogue entry
\subsection{Definitions and mathematical statement}
Let $G$ be finitely generated, and identify $H^1(G;\R)$ with $\operatorname{Hom}(G,\R)$. Its character sphere is $S(G)=(H^1(G;\R)\setminus\{0\})/\R_{>0}$. For a nonzero character $\chi$, the Bieri--Neumann--Strebel invariant contains $[\chi]$ when the full subgraph on vertices $\{g\in G:\chi(g)\geq0\}$ of a Cayley graph for a finite generating set is connected. To compare with a subset of cohomology, use its conical lift \[\widetilde\Sigma^1(G)=\{\chi\ne0:[\chi]\in\Sigma^1(G)\}\subseteq H^1(G;\R).\] Define the real first resonance variety by \[\mathcal R_1^1(G;\R)=\{a\in H^1(G;\R):\dim_\R H^1(H^*(G;\R),a\smile-)\geq1\}.\] The differential is multiplication by $a$ in the cohomology ring; its square is zero over $\R$. Equivalently, at degree one this cohomology is the kernel of $a\smile-:H^1\to H^2$ modulo the image of $a\smile-:H^0\to H^1$. If the source uses complex resonance, its real locus must be used for the displayed comparison. The source asks for the relationship between these invariants and, specifically, for conditions under which \[\widetilde\Sigma^1(G)\cup\bigl(-\widetilde\Sigma^1(G)\bigr)=H^1(G;\R)\setminus\mathcal R_1^1(G;\R).\] This retains the source's union, not an intersection. The conical lift resolves its suppressed sphere-versus-vector-space convention; equivalently one can compare nonzero rays after projectivizing both sides. The zero class needs a convention when $H^1(G;\R)=0$: with the displayed resonance definition the resonance set is empty, so the right side contains zero while the conical left side is empty. The literal equality and its nonzero-ray version then differ. No unprovided formality, finite-presentation, or finiteness-of-classifying-space assumption is imposed. The supplied 2026-08-21 triage labels the question open and reports Papadima--Suciu bounds involving BNSR invariants and homology jumping loci, with resonance descriptions under additional hypotheses, without a general characterization for all finitely generated groups.
\subsection{Short English statement}
Determine when the union of the BNS character cone and its negative is exactly the complement of first resonance. The sphere/cone convention, real resonance locus, and treatment of the zero character must be made explicit.
\subsection{Sources}
\begin{itemize}\item Local UnsolvedMath registry, permanent ID 30000825, source problem catalogue entry, \emph{BNS Invariants and First Resonance Varieties}, proposed year 2007; no individual proposer supplied.\item Supplied citation: \emph{Mini-Workshop: Topology of closed one-forms and Cohomology Jumping Loci} (2008), Oberwolfach Reports, DOI 10.4171/owr/2007/40, pp.\ 2321--2360: \url{https://doi.org/10.4171/owr/2007/40}. Supplied literature label: original Oberwolfach problem.\item Supplied literature label: Papadima--Suciu, \url{https://arxiv.org/abs/0812.2660}.\end{itemize}
\end{document}
